\documentclass[../main.tex]{subfiles}
\begin{document}

\chapter{メモリ}
\lessoninfo{
  video={Lesson 15，動画 1--14},
  textbook={H\&P \cite{hennessy2017} §2.2},
  keywords={SRAM，DRAM，メモリセル，ワード線，ビット線，センスアンプ，破壊読み出し，トレンチセル，リフレッシュ，行デコーダ，RAS と CAS，ファストページモード，フロントサイドバス，ノースブリッジ},
}

第12〜14章では，主記憶をキャッシュの背後に隠れた大きく遅い装置として
扱ってきた．また第1章では，プロセッサとメモリの速度の差，すなわちメモリ
ウォールが広がり続けることを見た．この章ではメモリそのものの内部を見る．
まず，メモリを作る2つの技術---キャッシュ用の静的 RAM と主記憶用の動的
RAM---を，1ビットを蓄える回路まで掘り下げて述べる．続いて，DRAM セルがどの
ようにチップ上に構成されアクセスされるか，そしてチップがどのようにプロセッサ
に接続されるかを示す．これらを合わせると，主記憶が遅い理由と，主記憶が
キャッシュのより速い技術で作られていない理由が分かる．

\section{メモリ技術}
\label{sec:l15-technologies}

キャッシュも主記憶も，ランダムアクセスメモリである．\source{Lesson 15，
動画 1--3；H\&P §2.2．}

\begin{definition}[ランダムアクセス]
  どの位置もそのアドレスによって直接，しかもどの位置についてもほぼ同じ時間
  で読み書きできるとき，そのメモリは
  \term[らんだむあくせす]{ランダムアクセス}[random access]を提供するという．
\end{definition}

その反対が順次アクセス（第16章）である．磁気テープのような媒体では，現在の
位置と要求された位置の間にあるデータすべてを通り過ぎるまで媒体を動かさなけれ
ばならない．ランダムアクセスメモリ（RAM）は2つの技術で作られ，その名前は，
電源が入っている間に，記憶されたビットに何が起こるかに由来する．

\begin{definition}[SRAM と DRAM]
  \term{SRAM}[static random-access memory]では，記憶されたビットは電源が
  入っている限りその値を保つ．\term{DRAM}[dynamic random-access memory]では，
  記憶されたビットは，電源が入っていても，周期的に書き直さない限り短時間で
  値を失う．
\end{definition}

これらの名前は，電源がないときに何が起こるかについては何も言っていない．
どちらの技術も揮発性であり，電源を切れば内容をすべて失う．

2つの技術は1ビットを蓄える回路が異なり，この違いが速度とコストを決める．
SRAM は1ビット当たり複数のトランジスタを使い，通常の設計では6個である
（\cref{sec:l15-sram}）．DRAM は1個のトランジスタしか使わない
（\cref{sec:l15-dram}）．したがって SRAM は1ビット当たりのチップ面積が
大きくビット当たりの価格も高いが，速い．そのセルがビット線を能動的に駆動
するのに対し，DRAM セルはわずかな電荷を保持するだけで，その電荷は検出し，
書き戻し，リフレッシュしなければならないからである．DRAM は同じ面積により
多くのビットを詰め込め，ビット当たりの価格も安いが，遅い．速くなければ
ならず小さくてよいキャッシュは SRAM で作られ，大きく安くなければならない
主記憶は DRAM で作られる．\margin{この区別は絶対ではない．IBM の POWER7
（2010）は，プロセッサダイ上の DRAM セル（トレンチ型の eDRAM）で 32~MB の
L3 キャッシュを作った．}

\section{SRAM セル}
\label{sec:l15-sram}

メモリはビットを，同一の回路からなる2次元配列に蓄える．\source{Lesson 15，
動画 4；H\&P §2.2．}1ビットを蓄える回路を
\term[めもりせる]{メモリセル}[memory cell]という．セルは行と列に並べられ，
配列を横切る2組の配線を通じてアクセスされる．

\begin{definition}[ビット線とワード線]
  \term[びっとせん]{ビット線}[bit line]はセルの列に沿って走り，その列のセル
  から読まれる値，あるいはそのセルに書かれる値を運ぶ．
  \term[わーどせん]{ワード線}[word line]はセルの行に沿って走り，活性化される
  と，その行のすべてのセルをそれぞれの列のビット線に接続する．
\end{definition}

SRAM セル（\cref{fig:l15-cells}a）では，ビットは輪をなすように接続された
2つのインバータが保持する．各インバータの出力が他方の入力を駆動する．
最初のインバータが 1 を出力すれば2番目は 0 を出力し，それが今度は最初の
インバータを 1 に保つ．この輪は2つの状態のどちらでも安定で，インバータに
電源が入っている限りその状態を保つ．これが SRAM を静的たらしめている．
ゲートがワード線に接続されたアクセストランジスタが，輪の一方のノードを
ビット線に接続する．

\begin{figure}[htbp]
  \centering
  % One memory bit.  (a) 6T SRAM cell: two cross-coupled inverters hold the
% bit; two access transistors, gated by the word line, connect the two
% nodes of the loop to the bit lines B and not-B.  (b) 1T DRAM cell: one
% access transistor connects the bit line to a storage capacitor.
\begin{tikzpicture}[x=1cm,y=1cm, font=\footnotesize\sffamily, line cap=round,
  dot/.style={circle, fill, inner sep=0pt, minimum size=3.2pt},
  lab/.style={font=\scriptsize\sffamily, text=ln-muted}]
  % \lnnmos{x}{y}: horizontal access transistor, channel centred at x, leads
  % at height y, gate wire going up to the word line at y=2.4
  \def\lnnmos#1#2{%
    \draw (#1-0.3,#2) -- (#1-0.3,#2+0.15) -- (#1+0.3,#2+0.15) -- (#1+0.3,#2);
    \draw[very thick] (#1-0.3,#2+0.3) -- (#1+0.3,#2+0.3);
    \draw (#1,#2+0.3) -- (#1,2.4) node[dot] {};
  }
  % ================================================= (a) SRAM cell
  \draw[ln-accent, thick] (0,2.4) -- (6.0,2.4);
  \node[lab, anchor=south] at (3.0,2.42) {\iflangja{ワード線}{word line}};
  \draw[ln-accent, thick] (0.4,2.8) -- (0.4,-1.0) node[below, text=black] {$B$};
  \draw[ln-accent, thick] (5.6,2.8) -- (5.6,-1.0) node[below, text=black] {$\overline{B}$};
  % access transistors
  \lnnmos{1.25}{1.4}
  \draw (0.4,1.4) node[dot] {} -- (0.95,1.4);
  \draw (1.55,1.4) -- (2.2,1.4);
  \lnnmos{4.75}{1.4}
  \draw (3.8,1.4) -- (4.45,1.4);
  \draw (5.05,1.4) -- (5.6,1.4) node[dot] {};
  % storage nodes
  \draw (2.2,1.4) -- (2.2,-0.2) -- (2.8,-0.2);
  \draw (3.8,1.4) -- (3.8,-0.2) -- (3.25,-0.2);
  % upper inverter (left to right)
  \draw[fill=ln-box] (2.75,0.45) -- (2.75,0.95) -- (3.2,0.7) -- cycle;
  \draw[fill=white] (3.26,0.7) circle (0.06);
  \draw (2.2,0.7) node[dot] {} -- (2.75,0.7);
  \draw (3.32,0.7) -- (3.8,0.7) node[dot] {};
  % lower inverter (right to left)
  \draw[fill=ln-box] (3.25,-0.45) -- (3.25,0.05) -- (2.8,-0.2) -- cycle;
  \draw[fill=white] (2.74,-0.2) circle (0.06);
  \draw (2.68,-0.2) -- (2.2,-0.2);
  \node[lab] at (3.0,1.05) {\iflangja{インバータ対}{inverter loop}};
  \node[anchor=north] at (3.0,-1.45) {\iflangja{(a) SRAM セル（6 トランジスタ）}{(a) SRAM cell (6 transistors)}};
  % ================================================= (b) DRAM cell
  \draw[ln-accent, thick] (7.2,2.4) -- (10.6,2.4);
  \node[lab, anchor=south] at (9.7,2.42) {\iflangja{ワード線}{word line}};
  \draw[ln-accent, thick] (7.8,2.8) -- (7.8,-1.0) node[below, text=black, lab] {\iflangja{ビット線}{bit line}};
  \lnnmos{8.65}{1.4}
  \draw (7.8,1.4) node[dot] {} -- (8.35,1.4);
  \draw (8.95,1.4) -- (9.6,1.4) -- (9.6,0.55);
  % capacitor and ground
  \draw[very thick] (9.3,0.55) -- (9.9,0.55);
  \draw[very thick] (9.3,0.38) -- (9.9,0.38);
  \draw (9.6,0.38) -- (9.6,-0.2);
  \draw (9.38,-0.2) -- (9.82,-0.2);
  \draw (9.46,-0.3) -- (9.74,-0.3);
  \draw (9.54,-0.4) -- (9.66,-0.4);
  \node[lab, anchor=west] at (9.95,0.47) {\iflangja{キャパシタ}{capacitor}};
  \node[anchor=north] at (8.9,-1.45) {\iflangja{(b) DRAM セル（1 トランジスタ）}{(b) DRAM cell (1 transistor)}};
\end{tikzpicture}

  \caption{1ビットの記憶．(a)~SRAM セルは2つのインバータの輪にビットを
    保持する．輪のノードは，2つのアクセストランジスタによって相補的な
    ビット線 $B$ と $\overline{B}$ に接続される．(b)~DRAM セルは，1個の
    アクセストランジスタの奥にあるキャパシタの電荷としてビットを保持する．}
  \label{fig:l15-cells}
\end{figure}

\begin{description}
  \item[書き込み] ビット線を新しい値に強く駆動し，ワード線を活性化する．
    ビット線は，ノードを駆動している小さなインバータを押し負かす．輪は
    ビット線に一致する状態に落ち着き，ワード線を非活性化した後もその状態を
    保つ．
  \item[読み出し] ワード線を活性化すると，輪がビット線を駆動する．セルは
    小さくなければならないのでインバータのトランジスタは小さいが，ビット線は
    列のすべてのセルに接続された長い配線である．1個のセルはビット線の電圧を
    ゆっくりとしか変えられない．
\end{description}

読み出しを速めるために，ビットの反転値を保持する輪のもう一方のノードを，
2つ目のアクセストランジスタを通じて2本目のビット線 $\overline{B}$ に
接続する．ワード線が活性化されると，セルは2本のビット線を互いに逆向きに
引っ張り，値は両者を比較することで得られる．

\begin{definition}[センスアンプ]
  \term[せんすあんぷ]{センスアンプ}[sense amplifier]は，ビット線上の小さな
  電圧変化，あるいは1対のビット線の間の小さな差を検出し，それを完全な
  論理 0 または 1 に増幅する．
\end{definition}

相補的な対があれば，センスアンプは2本のビット線を比較するだけでよい．セルが
2本を逆向きに動かすので，ごくわずかな振れの後には差の符号が明らかであり，
基準電圧も要らない．1本のビット線だけであれば，その値を確実に判定できるまで
にはるかに大きく動かなければならない．この対は書き込みにも役立つ．値と
その反転値が輪の両側に同時に働くからである．こうしてできるセルは，各インバータ
に2個とアクセストランジスタ2個の計6個のトランジスタをもち，ワード線1本，
ビット線2本，そしてインバータの電源接続を必要とする．\margin{第12〜14章の
キャッシュの配列はすべてこのセルでできている．32~KB の L1 キャッシュは
データ部だけで \num{262144} ビット，約 160 万個のトランジスタである．}

\section{DRAM セル}
\label{sec:l15-dram}

DRAM セルはビットをキャパシタ上の電荷として蓄える
（\cref{fig:l15-cells}b）．\source{Lesson 15，動画 5--6；H\&P §2.2；
\cite{dennard1968}．}ワード線が制御する1個のアクセストランジスタが，
キャパシタをビット線に接続する．充電されたキャパシタが 1 を，放電された
キャパシタが 0 を表す．書き込むにはビット線をその値にしてワード線を活性化
する．するとキャパシタはトランジスタを通じて充電または放電される．ワード線を
非活性化するとトランジスタが遮断し，電荷を孤立させる．

しかし，遮断したトランジスタは完全なスイッチではない．\caution{SRAM の
「静的」は電力ではなく保持のことである．SRAM セルも漏れ電流をもち，プロセッサ
の静的電力の大きな部分をキャッシュが占める．}静的電力（第1章）を生む
のと同じ不完全さによってわずかな電流が流れ続け，記憶された 1 の電荷は少しずつ
漏れ出して，やがてそのセルは 0 と区別できなくなる．これが
DRAM を動的たらしめている．ビットは，電荷が失われすぎる前に書き直された
ときだけ生き延びる．

\begin{definition}[リフレッシュ]
  DRAM セルのビットを読み出し，電荷が漏れ出してしまう前に完全な電荷で書き
  戻すことを\term[りふれっしゅ]{リフレッシュ}[refresh]という．どのセルも，
  そのデータが使われているかどうかに関わらず，周期的にリフレッシュしなければ
  ならない．
\end{definition}

DRAM セルの読み出しも繊細である．まずビット線を 0 と 1 の中間の電圧にし，
それからワード線を活性化する．キャパシタはビット線と電荷を分け合い，セルが
1 を保持していればビット線の電圧はわずかに上がり，0 を保持していればわずかに
下がる．セルのキャパシタは長いビット線の容量に比べて極めて小さいので，この
変化は非常に小さく，検出にはセンスアンプが必要である．電荷を分け合った後，
キャパシタは中間の電圧の近くに取り残される．すなわち，記憶されていた値は
失われている．

\begin{definition}[破壊読み出し]
  \term[はかいよみだし]{破壊読み出し}[destructive read]とは，記憶された値を
  壊してしまう読み出しのことである．したがって，DRAM セルの読み出しの後には
  必ず，読み出した値をそのセルに書き戻す動作が続く．
\end{definition}

SRAM の読み出しは破壊的ではない．インバータの輪がノードを駆動し続け，値を
自力で復元するからである．

キャパシタは面積の問題も引き起こす．チップ表面に平らに作った場合，確実に検出
できるだけの大きさのキャパシタはトランジスタよりはるかに広い面積を占め，
1トランジスタのセルがもつ密度上の利点の多くが失われてしまう．

\begin{definition}[トレンチセル]
  \term[とれんちせる]{トレンチセル}[trench cell]とは，アクセストランジスタの
  隣のシリコンに垂直に掘った，深く狭い溝（トレンチ）の中にキャパシタを形成
  した DRAM セルである（\cref{fig:l15-trench}）．
\end{definition}

\begin{marginfigure}
  % Cross-section of a DRAM trench cell (schematic): the access transistor
% sits at the surface; the storage capacitor is a deep, narrow trench
% etched into the silicon next to it.
\resizebox{\linewidth}{!}{%
\begin{tikzpicture}[x=1cm,y=1cm, font=\footnotesize\sffamily,
  lab/.style={font=\footnotesize\sffamily, text=ln-muted}]
  % substrate
  \fill[ln-box] (0,0) rectangle (4.2,-3.2);
  \draw[ln-rule] (0,0) -- (4.2,0);
  % source/drain regions
  \draw[fill=ln-accent!30] (0.5,0) rectangle (1.3,-0.35);
  \draw[fill=ln-accent!30] (2.1,0) rectangle (3.5,-0.35);
  % gate (word line)
  \draw[fill=ln-accent] (1.3,0.08) rectangle (2.1,0.36);
  \draw (1.7,0.36) -- (1.7,0.9);
  \node[lab, anchor=south] at (1.7,0.9) {\iflangja{ワード線}{word line}};
  % bit line contact
  \draw (0.9,0) -- (0.9,0.6) -- (0.2,0.6);
  \node[lab, anchor=south west] at (0.0,0.62) {\iflangja{ビット線}{bit line}};
  % trench capacitor: storage electrode joined to the drain region, lined
  % with a thin dielectric
  \fill[ln-accent!30] (2.68,-0.3) rectangle (3.32,-2.82);
  \draw[ln-muted, line width=1.2pt] (2.64,-0.35) -- (2.64,-2.86) -- (3.36,-2.86) -- (3.36,-0.35);
  \node[rotate=-90, font=\footnotesize\sffamily] at (3.0,-1.6) {\iflangja{容量}{capacitor}};
  \node[lab, anchor=west] at (3.5,-2.55) {Si};
\end{tikzpicture}}

  \caption{トレンチセルの断面（模式図）．}
  \label{fig:l15-trench}
\end{marginfigure}

トレンチによって，キャパシタはチップ表面をほとんど使わずに広い電極面積を
得る．もう1つのよく使われる設計は，現在の DRAM の大半が採用しているもので，
キャパシタをアクセストランジスタの上に積み重ねる．いずれにせよ，セルが占める
面積はトランジスタよりわずかに大きい程度である．また，SRAM セルより配線も
少なくて済む．ビット線は2本ではなく1本であり，インバータのための電源接続も
要らない．\cref{tab:l15-sram-dram} に2つのセルの比較をまとめる．

\begin{table}[htbp]
  \centering\small
  \caption{SRAM と DRAM の比較．}
  \label{tab:l15-sram-dram}
  \begin{tabular}{@{}lll@{}}
    \toprule
                            & SRAM             & DRAM \\
    \midrule
    記憶素子                & インバータの輪   & キャパシタ \\
    1ビット当たりトランジスタ数 & 6            & 1 \\
    セル当たりビット線数    & 2                & 1 \\
    電源投入中の値の保持    & 保つ             & リフレッシュすれば保つ \\
    読み出し                & 非破壊           & 破壊的，書き戻しが続く \\
    速度                    & 速い             & 遅い \\
    1ビット当たり面積とコスト & 大きい          & 小さい \\
    用途                    & キャッシュ       & 主記憶 \\
    \bottomrule
  \end{tabular}
\end{table}

\begin{keypoint}
  DRAM はビットを1個のトランジスタの奥のわずかな電荷として蓄える．そのため
  密度が高く安価だが遅い．読み出しは必ず電荷を壊すので書き戻しが続かなければ
  ならず，漏れ電流が周期的なリフレッシュを強いるからである．SRAM は1ビット
  当たり6個のトランジスタという代償で，そのどちらも避けている．
\end{keypoint}

\section{メモリチップの構成}
\label{sec:l15-chip}

メモリチップのセルは $2^r$ 行 $2^c$ 列の配列をなし（\cref{fig:l15-chip}），
$r+c$ ビットのアドレスは $r$ ビットの行アドレスと $c$ ビットの列アドレスに
分割される．例えば，アドレスの上位ビットと下位ビットである．\source{Lesson
15，動画 7；H\&P §2.2．}DRAM チップの1ビットの読み出しは4つの段階で進む．
\begin{enumerate}
  \item \term[ぎょうでこーだ]{行デコーダ}[row decoder]が行アドレスを，
    $2^r$ 本のワード線のうちちょうど1本の活性化に変換する．
  \item 選ばれた行のすべてのセルがそれぞれのビット線に接続される．ビット線
    1本につき1個のセンスアンプが値を検出し，行全体が行ラッチに保持される．
  \item \term[れつまるちぷれくさ]{列マルチプレクサ}[column multiplexer]が
    列アドレスを使ってラッチの1ビットを選び，データ出力へ送り出す．
  \item 読み出しは行のすべてのセルにとって破壊的だったので，ラッチの内容を
    ビット線を通じて行に書き戻す．
\end{enumerate}

\begin{figure}[htbp]
  \centering
  % Organization of a DRAM chip: the row decoder activates one word line
% (row address or refresh counter), the sense amplifiers read the whole row
% into the row latch, the column multiplexer selects one bit, and the latch
% is written back to the row over the bit lines.
\begin{tikzpicture}[x=1cm,y=1cm, font=\footnotesize\sffamily, >={Stealth[length=1.8mm]},
  lab/.style={font=\scriptsize\sffamily, text=ln-muted},
  blk/.style={draw, fill=ln-box}]
  % ---- cell array: one cell at every crossing of a word line and a bit line
  \fill[ln-accent!15] (3.3,2.55) rectangle (9.3,3.05);
  \foreach \i in {0,...,5} {
    \draw[ln-accent] (3.3,{4.0-0.6*\i}) -- (9.3,{4.0-0.6*\i});
  }
  \foreach \j in {0,...,7} {
    \draw[ln-muted] ({3.9+0.7*\j},4.4) -- ({3.9+0.7*\j},0.3);
    \foreach \i in {0,...,5} {
      \draw[fill=white] ({3.9+0.7*\j-0.11},{4.0-0.6*\i-0.11}) rectangle ({3.9+0.7*\j+0.11},{4.0-0.6*\i+0.11});
    }
  }
  \node[lab, anchor=south] at (6.35,4.4) {\iflangja{メモリセルの配列}{array of memory cells}};
  \node[lab, anchor=west] at (9.35,4.0) {\iflangja{ワード線}{word line}};
  \node[lab, anchor=south] at (8.8,4.4) {\iflangja{ビット線}{bit line}};
  % ---- row decoder with address / refresh counter selection
  \draw[blk] (2.6,0.6) rectangle (3.3,4.4);
  \node[rotate=90] at (2.95,2.5) {\iflangja{行デコーダ}{row decoder}};
  \draw[fill=white] (1.7,1.75) -- (2.2,2.05) -- (2.2,2.95) -- (1.7,3.25) -- cycle;
  \draw[->] (2.2,2.5) -- (2.6,2.5);
  \draw[->] (0,2.9) node[above right, inner sep=1pt] {\iflangja{行アドレス}{row address}} -- (1.7,2.9);
  \node[blk, align=center, font=\scriptsize\sffamily, inner sep=2pt] (cnt) at (0.65,1.6)
    {\iflangja{リフレッシュ}{refresh}\\\iflangja{カウンタ}{counter}};
  \draw[->] (cnt.east) -- (1.45,1.6) -- (1.45,2.1) -- (1.7,2.1);
  % ---- sense amplifiers and row latch
  \draw[blk] (3.6,-0.2) rectangle (9.2,0.3);
  \node at (6.4,0.05) {\iflangja{センスアンプ}{sense amplifiers}};
  \draw[blk] (3.6,-1.3) rectangle (9.2,-0.8);
  \node at (6.4,-1.05) {\iflangja{行ラッチ}{row latch}};
  % read into the latch and write back over the same bit lines
  \foreach \j in {0,...,7} {
    \draw[<->, >={Stealth[length=1.5mm]}] ({3.9+0.7*\j},-0.2) -- ({3.9+0.7*\j},-0.8);
  }
  \node[lab, anchor=west, align=left] at (9.25,-0.5)
    {\iflangja{読み出し／}{read /}\\\iflangja{書き戻し}{write back}};
  % ---- column multiplexer
  \foreach \j in {0,...,7} {
    \draw ({3.9+0.7*\j},-1.3) -- ({3.9+0.7*\j},-1.6);
  }
  \draw[fill=white] (3.7,-1.6) -- (9.1,-1.6) -- (7.15,-2.3) -- (5.65,-2.3) -- cycle;
  \node at (6.4,-1.92) {\iflangja{列マルチプレクサ}{column multiplexer}};
  \draw[->] (0.8,-1.95) node[above right, inner sep=1pt] {\iflangja{列アドレス}{column address}} -- (4.68,-1.95);
  \draw[->] (6.4,-2.3) -- (6.4,-2.8) node[below] {\iflangja{データ}{data}};
\end{tikzpicture}

  \caption{DRAM チップの構成．行デコーダは，行アドレスまたはリフレッシュ
    カウンタ（\cref{sec:l15-refresh}）が選ぶワード線を活性化する．センス
    アンプが行全体を行ラッチに読み出し，列マルチプレクサが1ビットを選び，
    ラッチが行に書き戻される．}
  \label{fig:l15-chip}
\end{figure}

書き込みも同じ経路を使う．上と同じように行をラッチに読み出し，列アドレスが
選ぶビットをラッチ上で書き込むべき値に置き換え，書き戻しによって変更された
行を格納する．このように，DRAM チップにおける書き込みとは，書き戻す前に
ラッチを変更する読み出しである．\margin{実際のデバイスの幅は1ビットではなく
4，8，16 ビットである．8 ビット幅の DDR4 チップは1行に 1~KB をもち，ランクの
8個のチップは同時に行を開く（JEDEC DDR4）．}

\begin{note}
  \label{note:l15-model}
  行を閉じるときに書き戻される行ラッチは，この章を通じて用いる単純化した
  モデルである．実際の DRAM チップでは，センスアンプ自身が行を保持する．
  ワード線が活性のままの間にセンスアンプがビット線を完全なレベルまで駆動し，
  それによって行のセルが復元される．また，ラッチに書き込まれたビットは
  ただちにそのセルに届く．行を閉じる操作はワード線を落とし，ビット線を次の
  読み出しに備えさせる．
\end{note}

アドレスの2つの部分にかかる時間は大きく異なる．行を選ぶには，デコードし，
長いワード線を駆動し，長いビット線上のわずかな電荷をセンスアンプが検出する
のを待たなければならない．列を選ぶのは，すでにラッチにあるビットの間で
マルチプレクサを切り替えるだけである．\cref{sec:l15-ras-cas}と
\cref{sec:l15-fpm}はこの差を利用する．\margin{データシートではラッチされた行を
行バッファと呼び，すでに開いている行に当たるアクセスを「ページヒット」という
（\cref{sec:l15-fpm}）．}

\begin{note}
  SRAM チップも同じように構成され，行デコーダ，センスアンプ，列マルチ
  プレクサをもつ．ただし，その読み出しは破壊的ではないので，読み出しの後に
  書き戻しは続かず，またそのセルにリフレッシュも要らない．
\end{note}

\section{リフレッシュ}
\label{sec:l15-refresh}

DRAM チップのどの行も，長さ $T$ の区間ごとに少なくとも1回はリフレッシュ
しなければならない．$T$ は普通，数十ミリ秒である．\source{Lesson 15，
動画 8--9；H\&P §2.2．}リフレッシュは通常の読み出し経路を使う．すなわち，出力を
使わずにラッチを書き戻す読み出し（\cref{sec:l15-chip}）である．\margin{JEDEC の
DDR4 は \SI{64}{\milli\second} の窓に \num{8192} 回のリフレッシュコマンドを
収め，\SI{7.8}{\micro\second} ごとに1回発行する．\SI{85}{\celsius} を超えると
間隔は半分になる．}

プログラムによる読み出しは，触れた行をリフレッシュする．しかし，そのような
読み出しは当てにできない．長い間，あるいはまったく読まれない行があるから
である．そこで，すべての行を系統的にリフレッシュする．リフレッシュカウンタが
行番号を保持する．周期的に，行アドレスの代わりにカウンタの値が行デコーダに
与えられ（\cref{fig:l15-chip}），その行が読み出されて書き戻され，カウンタは
次の行へ進み，最後の行の後は先頭に戻る．行数を $R$ とし，そのそれぞれを
区間 $T$ ごとに1回リフレッシュしなければならないとすると，チップは，次の長さ
の区間ごとに少なくとも1回，1つの行をリフレッシュしてカウンタを進めなければ
ならない．
\begin{equation}
  t_{\text{row}} = \frac{T}{R} .
  \label{eq:l15-refresh}
\end{equation}

ある行をリフレッシュしている間，チップは読み出しにも書き込みにも応じられず，
その間に到着したアクセスは待たなければならない．こうしてリフレッシュは
DRAM が遅い理由をもう1つ増やし，そのコストは行数とともに増える．\tip{リフレッシュ
のオーバーヘッドは $t/(T/R) = Rt/T$ である．$t$ は1行のリフレッシュがチップを
占有する時間であり，$T$ の中でどう散らすかは関係しない．}

\begin{example}
  \label{ex:l15-refresh}
  ある DRAM チップは \num{8192} 行をもち，そのセルは \SI{64}{\milli\second}
  ごとにリフレッシュしなければならない．\cref{eq:l15-refresh} により，チップ
  は $\SI{64}{\milli\second}/\num{8192} \approx \SI{7.8}{\micro\second}$
  ごとに1行をリフレッシュしなければならない．すなわち毎秒 \num{128000} 行で
  ある．1行のリフレッシュに \SI{50}{\nano\second} かかるとすると，チップは
  毎秒のうち $\num{128000} \times \SI{50}{\nano\second} =
  \SI{6.4}{\milli\second}$，すなわち時間の \SI{0.64}{\percent} を
  リフレッシュに費やす．行数が4倍で $T$ が同じチップでは
  \SI{2.56}{\percent} となる．
\end{example}

\section{アドレスの多重化}
\label{sec:l15-ras-cas}

$2^{r+c}$ 個のセルをもつチップには $r+c$ ビットのアドレスが必要であり，
並列に送るビットはどれもパッケージのピン1本とチップへの配線1本を消費する．
\source{Lesson 15，動画 10；H\&P §2.2．}しかし，行アドレスと列アドレスが
同時に必要になるわけではない．行アドレスは，遅い行の選択を始めるために先に
必要である．列アドレスは，行がラッチに入ってから初めて必要になる．そこで
DRAM チップは両方の部分を同じピンから順に受け取り，2本の制御信号が，どちらの
部分が来ているかをチップに伝える．

\begin{definition}[RAS と CAS]
  \term[ぎょうあどれすすとろーぶ]{行アドレスストローブ}[row address strobe]%
  （RAS）がアサートされると，チップはアドレスピン上の値を行アドレスとして
  取り込み，ワード線を活性化して行をラッチに読み出す．
  \term[れつあどれすすとろーぶ]{列アドレスストローブ}[column address strobe]%
  \index{CAS|see{列アドレスストローブ}}（CAS）がアサートされると，チップは
  アドレスピン上の値を列アドレスとして取り込み，ラッチのそのビットを選ぶ．
\end{definition}

読み出しでは，選ばれたビットがデータ出力に現れる．CAS と同時にライト
イネーブル信号がアサートされていれば，代わりに，ラッチ中のそのビットが
データ入力の値に置き換えられる．RAS をデアサートするとアクセスが終わり，
（\cref{note:l15-model} のモデルでは）ラッチが書き戻されて行が閉じられる．
\cref{fig:l15-timing}a は，この種のアクセスを2回示している．\margin{実際の
チップでは RAS と CAS は負論理である．\cref{fig:l15-timing} では
アサートされた信号を高レベルで描いている．}

\begin{figure}[htbp]
  \centering
  % RAS/CAS timing on a DRAM chip with multiplexed address pins (asserted
% signals drawn high).  (a) Two separate accesses to the same row: each
% opens the row with RAS, selects a column with CAS and closes the row.
% (b) Fast page mode: the row is opened once and three columns are
% accessed before RAS is deasserted and the row is written back.
\begin{tikzpicture}[x=1cm,y=1cm, font=\footnotesize\sffamily,
  lab/.style={anchor=east, font=\footnotesize\sffamily},
  ann/.style={font=\scriptsize\sffamily, text=ln-muted}]
  % \lnbus{x0}{x1}{y}{text}: valid value on the address pins
  \def\lnbus#1#2#3#4{%
    \draw[fill=ln-accent!15] (#1,#3) -- (#1+0.1,#3+0.2) -- (#2-0.1,#3+0.2) -- (#2,#3)
      -- (#2-0.1,#3-0.2) -- (#1+0.1,#3-0.2) -- cycle;
    \node[font=\scriptsize\sffamily] at ({(#1+#2)/2},#3) {#4};
  }
  % \lnidle{x0}{x1}{y}: no value on the address pins
  \def\lnidle#1#2#3{\draw[ln-muted] (#1,#3) -- (#2,#3);}
  % \lnsig{y}{list of x0/x1 pulses}: control signal, low outside the pulses
  \def\lnsig#1#2{%
    \draw (1.4,#1-0.2) \foreach \a/\b in {#2} { -- (\a,#1-0.2) -- (\a+0.08,#1+0.2)
      -- (\b-0.08,#1+0.2) -- (\b,#1-0.2) } -- (11.4,#1-0.2);
  }
  % ================================================= (a)
  \node[anchor=west] at (0,0.75) {\iflangja{(a) 同じ行への 2 回のアクセス}{(a) two accesses to the same row}};
  \node[lab] at (1.3,0) {\iflangja{アドレス}{address}};
  \lnidle{1.4}{2.0}{0}
  \lnbus{2.0}{3.4}{0}{\iflangja{行}{row}}
  \lnbus{3.4}{4.8}{0}{\iflangja{列}{column}}
  \lnidle{4.8}{6.6}{0}
  \lnbus{6.6}{8.0}{0}{\iflangja{行}{row}}
  \lnbus{8.0}{9.4}{0}{\iflangja{列}{column}}
  \lnidle{9.4}{11.4}{0}
  \node[lab] at (1.3,-0.6) {RAS};
  \lnsig{-0.6}{2.3/5.1, 6.9/9.7}
  \node[lab] at (1.3,-1.2) {CAS};
  \lnsig{-1.2}{3.7/4.7, 8.3/9.3}
  \node[ann, anchor=west] at (5.15,-0.6) {\iflangja{書き戻し}{write back}};
  \node[ann, anchor=west] at (9.75,-0.6) {\iflangja{書き戻し}{write back}};
  % ================================================= (b)
  \node[anchor=west] at (0,-2.05) {\iflangja{(b) ファストページモード}{(b) fast page mode}};
  \node[lab] at (1.3,-2.8) {\iflangja{アドレス}{address}};
  \lnidle{1.4}{2.0}{-2.8}
  \lnbus{2.0}{3.4}{-2.8}{\iflangja{行}{row}}
  \lnbus{3.4}{4.8}{-2.8}{\iflangja{列}{column}}
  \lnbus{4.8}{6.2}{-2.8}{\iflangja{列}{column}}
  \lnbus{6.2}{7.6}{-2.8}{\iflangja{列}{column}}
  \lnidle{7.6}{11.4}{-2.8}
  \node[lab] at (1.3,-3.4) {RAS};
  \lnsig{-3.4}{2.3/8.2}
  \node[lab] at (1.3,-4.0) {CAS};
  \lnsig{-4.0}{3.7/4.7, 5.1/6.1, 6.5/7.5}
  \node[ann, anchor=west] at (8.2,-3.4) {\iflangja{書き戻し}{write back}};
\end{tikzpicture}

  \caption{アドレスピンを多重化した DRAM チップの信号．(a)~各アクセスは
    RAS で行を開き，CAS で列を選び，行を閉じる．(b)~ファストページモード
    （\cref{sec:l15-fpm}）では，行を一度だけ開いて3つの列にアクセスし，
    RAS をデアサートしたときに行を書き戻す．}
  \label{fig:l15-timing}
\end{figure}

多重化により，アドレスピンは $r+c$ 本から $\max(r,c)$ 本に減る．行数と列数が
等しければ半分である．これによってパッケージは小さく安くなり，チップと，
それを駆動するメモリコントローラ（\cref{sec:l15-connection}）との間の配線も
減る．例えば，\num{16384} 行 \num{16384} 列に $2^{28}$ 個のセルをもつチップ
は，28 本ではなく 14 本のアドレスピンで済む．

\section{ファストページモード}
\label{sec:l15-fpm}

\cref{fig:l15-timing}a の流れでは，どのアクセスも行を開き，1つの列を使い，
また行を閉じる．次のアクセスが同じ行に対するものであってもそうである．
\source{Lesson 15，動画 11--12．}しかし，いったん行を開いてしまえば，その行
のビットはすべてラッチにあり，その行の別のビットに届くには新しい列アドレス
だけがあればよい．

\begin{definition}[ファストページモード]
  \term[ふぁすとぺーじもーど]{ファストページモード}[fast page mode]では，
  RAS で一度開いた行を開いたままにし，続けざまの CAS パルスで複数の列
  アドレスを与えて，そのたびに行の1ビットを読み書きする．行が書き戻されるの
  は，最後に RAS がデアサートされるときの1回だけである
  （\cref{fig:l15-timing}b）．
\end{definition}

\caution{ここでいう「ページ」はチップの行のことであり，仮想記憶のページ
（第13章）ではない．}

ファストページモードでは，高価な段階---行の選択とセンス，そして書き戻し---を
アクセスのまとまり全体につき1回だけ行い，各アクセスに要するのは速い列の選択
だけである．\sidenote{ファストページモードは一連の改良の最初のものだった．
EDO（extended data out）は次の列アドレスを与える間もデータを有効に保ち，
同期式 DRAM（SDRAM，最初の製品は 1992 年）はストローブをクロックとコマンドの
体系に置き換え，DDR はクロックの両エッジで転送するようにした．チップ内部の行と
列の構造は変わっていない．}空間的局所性（第12章）によって頻繁に起こる近いアドレスへの
アクセスは，普通はアドレスの下位ビットだけが異なるので，同じ行を共有する．
行をどれだけ長く開いたままにするかは，RAS と CAS を駆動するコントローラが
決める．後続のアクセスが同じ行を使う間は行を開いたままにでき，
別の行にアクセスするときや，どれかの行をリフレッシュするときには，その前に
行を閉じなければならない．したがって，節約できる量は，連続するアクセスの
うちどれだけが同じ行を共有するかによる．

\begin{example}
  \label{ex:l15-fpm}
  ある DRAM チップでは，行を開く（RAS，デコードとセンス）のに
  \SI{30}{\nano\second}，CAS で列を選ぶのに \SI{5}{\nano\second}，行を閉じる
  （書き戻し）のに \SI{10}{\nano\second} かかる．\cref{tab:l15-fpm} の順に
  5つのアクセスが到着する．表中の組は（行，列）を表す．ファストページモード
  がなければ，どのアクセスも自分の行を開いて閉じるので
  $5\times\SI{45}{\nano\second} = \SI{225}{\nano\second}$ かかる．ファスト
  ページモードでは行を開いて閉じるのが3回で済み，表の合計は
  \SI{145}{\nano\second} となる．行 22 へのアクセスが最後に来ていれば，
  行 7 への4つのアクセスが1回の開きを共有し，合計は \SI{105}{\nano\second}
  だったはずである．\caution{22-22-22 のようなモジュールのタイミングはナノ秒
  ではなくバスのクロック数である．DDR4-3200 では，行を開く，列を読む，行を
  閉じるのがそれぞれ \SI{13.75}{\nano\second} である．JEDEC のスピードビン．}
\end{example}

\begin{table}[htbp]
  \centering\small
  \caption{\cref{ex:l15-fpm} のアクセスをファストページモードで実行した
    場合．}
  \label{tab:l15-fpm}
  \begin{tabular}{@{}llr@{}}
    \toprule
    アクセス          & 動作                                    & 時間 (ns) \\
    \midrule
    読み出し (7, 12)  & 行 7 を開く，列 12 を CAS                & 35 \\
    読み出し (7, 13)  & 列 13 を CAS                            & 5 \\
    書き込み (7, 14)  & ライトイネーブルとともに列 14 を CAS      & 5 \\
    読み出し (22, 3)  & 行 7 を閉じる，行 22 を開く，列 3 を CAS  & 45 \\
    読み出し (7, 15)  & 行 22 を閉じる，行 7 を開く，列 15 を CAS & 45 \\
    ---               & 行 7 を閉じる                            & 10 \\
    \midrule
    合計              &                                         & 145 \\
    \bottomrule
  \end{tabular}
\end{table}

\section{DRAM とプロセッサの接続}
\label{sec:l15-connection}

プロセッサが主記憶を待つ時間は，DRAM チップのアクセス時間だけではない．
\source{Lesson 15，動画 13--14．}両者の間には他のチップとバスを通る経路が
あり，その伝統的な形を \cref{fig:l15-connection}a に示す．

\begin{definition}[フロントサイドバスとノースブリッジ]
  \term[ふろんとさいどばす]{フロントサイドバス}[front-side bus]%
  \index{FSB|see{フロントサイドバス}}（FSB）は，プロセッサを
  \term[のーすぶりっじ]{ノースブリッジ}[northbridge]に接続する．ノース
  ブリッジは，プロセッサを主記憶と，グラフィックスアダプタのような高速な
  デバイスと，そして遅い I/O デバイスが接続されるサウスブリッジ（第16章）と
  に結ぶチップである．
\end{definition}

ノースブリッジは\term[めもりこんとろーら]{メモリコントローラ}[memory controller]%
を内蔵する．メモリコントローラは，FSB を通じて届くメモリアクセスのそれぞれ
を DRAM チップの信号に翻訳する．すなわち，アドレスを行アドレスと列アドレスに
分割し，メモリバス上に RAS，CAS，ライトイネーブルを駆動する．したがって，
キャッシュミスによるメモリアクセスは，プロセッサから FSB を通って
ノースブリッジへ，そこからメモリバスを通って DRAM チップへと進み，データは
同じ道を戻る．この各段階がレイテンシを加える上に，FSB は，プロセッサと
ノースブリッジの先のデバイスとの間のそれ以外のトラフィックもすべて運ぶ．
\margin{規模の目安として，64 ビット幅で 1333~MT/s のフロントサイドバスは
すべてのトラフィックを合わせて約 10.7~GB/s を運ぶが，DDR4-3200 のチャネル
1本だけで 25.6~GB/s である．}

\begin{figure}[htbp]
  \centering
  % Path from the processor to DRAM.  (a) Traditional organization: the
% processor reaches the memory controller in the northbridge over the
% front-side bus; the northbridge drives the memory bus.  (b) Integrated
% memory controller: the memory bus is attached to the processor itself.
\begin{tikzpicture}[x=1cm,y=1cm, font=\footnotesize\sffamily, >={Stealth[length=1.8mm]},
  box/.style={draw, fill=ln-box, align=center, inner sep=3pt},
  mc/.style={draw, fill=ln-accent!15, align=center, inner sep=2pt, font=\scriptsize\sffamily},
  lab/.style={font=\scriptsize\sffamily, text=ln-muted}]
  % \lndimms{x}{y}: three DRAM modules centred at (x,y)
  \def\lndimms#1#2{%
    \foreach \k in {0,1,2} {
      \draw[fill=white] (#1-0.75+0.12*\k,#2-0.45+0.14*\k) rectangle (#1+0.55+0.12*\k,#2+0.25+0.14*\k);
    }
    \node[font=\scriptsize\sffamily] at (#1+0.04,#2+0.08) {DRAM};
  }
  % ================================================= (a)
  \node[anchor=west] at (0,1.9) {\iflangja{(a) ノースブリッジ経由}{(a) through the northbridge}};
  \node[box, minimum width=1.9cm, minimum height=1.0cm] (cpu) at (1.0,0) {\iflangja{プロセッサ}{processor}};
  \node[box, minimum width=2.3cm, minimum height=1.3cm] (nb) at (5.1,0) {};
  \node[anchor=north] at (5.1,0.6) {\iflangja{ノースブリッジ}{northbridge}};
  \node[mc] at (5.1,-0.25) {\iflangja{メモリコントローラ}{memory controller}};
  \lndimms{9.6}{0}
  \node[box, font=\scriptsize\sffamily] (gfx) at (5.1,1.5) {\iflangja{グラフィックス}{graphics}};
  \node[box, font=\scriptsize\sffamily] (sb) at (5.1,-1.55) {\iflangja{サウスブリッジ}{southbridge}};
  \draw[<->, very thick, ln-accent] (cpu) -- node[lab, above] {FSB} (nb);
  \draw[<->, very thick, ln-accent] (nb.east) -- node[lab, above] {\iflangja{メモリバス}{memory bus}} (8.8,0);
  \draw[<->] (nb) -- (gfx);
  \draw[<->] (nb) -- (sb);
  \draw[->] (sb.east) -- (7.0,-1.55) node[lab, anchor=west] {I/O};
  % ================================================= (b)
  \node[anchor=west] at (0,-2.55) {\iflangja{(b) メモリコントローラ統合}{(b) integrated memory controller}};
  \node[box, minimum width=4.3cm, minimum height=1.2cm, anchor=west] (cpu2) at (0.05,-3.5) {};
  \node[anchor=west] at (0.2,-3.5) {\iflangja{プロセッサ}{processor}};
  \node[mc, anchor=east] at (4.25,-3.5) {\iflangja{メモリコントローラ}{memory controller}};
  \lndimms{9.6}{-3.5}
  \draw[<->, very thick, ln-accent] (cpu2.east) -- node[lab, above] {\iflangja{メモリバス}{memory bus}} (8.8,-3.5);
\end{tikzpicture}

  \caption{プロセッサと DRAM の間の経路．(a)~ノースブリッジ内のメモリ
    コントローラには，フロントサイドバスを通じて到達する．(b)~メモリ
    コントローラをプロセッサに統合すると，メモリバスはプロセッサに直接
    接続される．}
  \label{fig:l15-connection}
\end{figure}

新しいプロセッサは，メモリコントローラをプロセッサチップ上に統合し，メモリ
バスをプロセッサに直接接続する（\cref{fig:l15-connection}b）．\sidenote{メモリ
コントローラをプロセッサダイに載せた最初の主流の x86 プロセッサは AMD の
Athlon~64（2003）であり，Intel は 2008 年の Nehalem で続いた．ノースブリッジの
残り---グラフィックスと PCI~Express のリンク---もプロセッサ側に移り，遅い
デバイスのためのチップセットが1つ残った．いまや各プロセッサチップが自分の
メモリをもち，他のチップのメモリへのアクセスは相互接続網を越える．第18章の
NUMA である．}こうすると
FSB とノースブリッジはメモリへの経路から消え，主記憶へのすべてのアクセスの
レイテンシが下がる．それでもなお，この章にまとめた理由により，DRAM の
アクセスはプロセッサよりはるかに遅く，第12〜14章のキャッシュは欠かせない
ままである．第16章では，さらに遅いが電源がなくてもデータを保つ記憶装置に
話を移す．

\begin{keypoint}[章のまとめ]
  RAM は，6個のトランジスタからなるセルがインバータの輪にビットを保持する
  SRAM として，あるいは1個のトランジスタからなるセル（例えばトレンチセル）
  がキャパシタ上の電荷としてビットを保持する DRAM として作られる．SRAM は
  速くて高価であり，キャッシュに使われる．DRAM は密度が高くて安価であり，
  主記憶に使われるが，その読み出しは破壊的で書き戻しが続き，漏れ電流が
  すべての行のリフレッシュを強いるので，$T/R$ の区間ごとに1行をリフレッシュ
  しなければならない．DRAM チップは，行デコーダで行を選び，センスアンプで
  それをラッチに読み出し，列マルチプレクサで1ビットを選ぶ．行アドレスと列
  アドレスはアドレスピンを共有し，RAS と CAS で区別される．そしてファスト
  ページモードは，1つの行への複数のアクセスを1回の開きで処理する．メモリ
  アクセスは，伝統的にはフロントサイドバスを通ってノースブリッジ内のメモリ
  コントローラへ，さらにメモリバスを通って進む．新しいプロセッサはメモリ
  コントローラを統合し，メモリバスを直接駆動する．
\end{keypoint}

\end{document}
